\section{Introduction}
\label{sec:introduction}

\texttt{diffmrc} is an end-to-end differentiable multireference
electronic structure stack.  It spans configuration interaction
(CASCI), complete active space self-consistent field (CASSCF),
state-averaged CASSCF, and second-order N-electron valence perturbation
theory (NEVPT2).  Every energy, nuclear gradient, and second derivative
the package produces comes out of the same code path by automatic
differentiation; there is no hand-derived derivative equation anywhere
in the source.

The stack is built on \href{https://pyscf.org/}{PySCF} and
\href{https://github.com/fishjojo/pyscfad}{pyscfad}, with
\href{https://github.com/jax-ml/jax}{JAX} providing the differentiable
array operations.  All arithmetic is double precision.  The package
does not replace PySCF: the Gaussian-basis integrals, the reference
SCF, and the benchmark comparisons all come from PySCF and pyscfad,
while the differentiable multireference chain is original to this
codebase.

\subsection{Who this tutorial is for}

You should be comfortable with the chemistry of active spaces and with
running a Python quantum-chemistry calculation from a script.  You do
not need prior exposure to JAX: the tutorial introduces the runtime
settings that matter (\texttt{jax\_enable\_x64}, the CPU platform pin,
and the thread budget) right where they are first needed.

\subsection{What this tutorial covers}

\begin{itemize}
  \item Chapter~\ref{sec:installation}: installing the package and
    setting up a working environment.
  \item Chapter~\ref{sec:first-calculation}: a complete CASSCF
    nuclear-gradient calculation for water, explained line by line.
  \item Chapter~\ref{sec:state-averaged}: state-averaged CASSCF with
    per-state NEVPT2, and how to read the root table it produces.
  \item Chapter~\ref{sec:active-space}: choosing an active space, with
    the frontier-window rules the benchmark work forced us to write
    down.
  \item Chapter~\ref{sec:convergence}: what the convergence
    diagnostics mean and how to act on them.
  \item Chapter~\ref{sec:benchmark-walkthrough}: a full benchmark
    walkthrough on a real target, end to end.
  \item Chapter~\ref{sec:limitations}: where the method family and the
    present implementation stop being reliable.
\end{itemize}

The companion \emph{Manual} carries what a tutorial should not: the
API reference, protocol tables, the design of the differentiable chain
and its second derivatives, the verification records, and the
development architecture.

\subsection{Project status}

This is a research codebase, released to accompany a methods paper and
to make its results reproducible.  It is not a finished library:
interfaces may change without notice, test coverage is partial, and
the numerical behaviour of the orbital-optimization and perturbation
layers depends on convergence settings that the code reports rather
than enforces.  Several limitations listed in
Chapter~\ref{sec:limitations} are properties of the method family and
will not be removed by further development of this package.  Use it
for research, read the convergence diagnostics it returns, and do not
treat it as a black box.
